\section{Lessons from the Replication Crisis}\label{sec:replication}

The failures described in the last two sections are not new. Over the past decade the social and behavioral sciences documented them in human analysts, and that work is the closest thing this review has to settled knowledge about what keeps an analysis honest. Its central lesson is that the problem lies in the structure of analysis rather than in the character of the analyst. The remedy is structural too.

\subsection{The Analyst Problem}

Many-analysts studies give the same data and question to independent teams and compare what comes back. When 29 teams asked whether soccer referees give more red cards to dark-skinned players, their estimates ranged from 0.89 to 2.93 in odds-ratio units; 20 teams found a significant positive effect and 9 did not. Neither prior beliefs, nor expertise, nor the peer-rated quality of the analyses ``readily explained'' the spread \citep{R9}. Among 70 neuroimaging teams, ``no two teams chose identical workflows,'' and on average 20\% reported a result that differed from the majority \citep{R10}. Across 1,261 models from 72 teams, coded research-design decisions explained 2.6\% of the variance in results and researcher characteristics at most 1.2\%, leaving 95.2\% unexplained \citep{R11}. The pattern recurs in ecology, with more than 170 analyst teams \citep{R16}, and in finance, where 164 teams testing the same hypotheses produced sizable ``nonstandard errors'' \citep{R17}. Gelman and Loken explain why honest intent does not help: when analysis is ``highly contingent on data,'' multiple comparisons arise without ``any conscious procedure'' of fishing for significance \citep{R12}. These studies measure disagreement, not error; they do not say which team was right. Replication projects measure the cost. Of 100 psychology studies, 97\% of originals and 36\% of replications were significant, with replication effects half as large \citep{R14}; high-powered replications of 21 experiments from \textit{Nature} and \textit{Science} found a significant same-direction effect for 13 (62\%) \citep{R18}.

\subsection{Structures That Work}

The reforms that followed commit decisions before results exist, or expose the decisions that were made. The evidence for pre-commitment is observational but consistent. Registered Reports in psychology reported positive results for their first hypothesis 44\% of the time, against 96\% in standard reports, although that study ``was not an experiment'' \citep{R13}. Among 15,992 test statistics from randomized trials in leading economics journals, registration alone made ``no meaningful difference,'' while a complete pre-analysis plan lowered the share of marginally significant results by about 10 percentage points \citep{R19}. What matters is the detail of the plan, not the timestamp. The reform can fail its own standard: a paper claiming that high replicability of new findings ``is achievable'' was retracted, in part for analyses selected ``with knowledge of the data'' \citep{R20}. Other tools make the forking paths visible. Multiverse tools generate and run every specified variant of an analysis \citep{R30}, the proposed m-value locates a reported result within that distribution \citep{R3}, and null checks rerun an analysis on data with the signal shuffled away \citep{R4}. Independent re-execution comes last. Of R files in replication packages on the Harvard Dataverse repository, 74\% failed on their first clean run \citep{R15}. In a preprint that re-ran 384 political-science papers with an agent-assisted pipeline, the reproducible share rose from 20.8\% before the journals adopted data-access and verification policies to 82.5\% after; the authors credit ``active verification---not merely archiving mandates'' \citep{R8}.

\subsection{Agents as Analysts at Machine Speed}

An agent can walk many forking paths in the time a person walks one, and the new agent studies suggest this changes the scale of the problem, not its nature. In 4,946 runs of LLM analysts, ``independent runs frequently reverse whether the hypothesis is judged supported''; the gap in support rates between a skeptical and a strongly confirmation-seeking persona ranged from 34 to 66 percentage points, and an AI auditor ``narrows, but does not eliminate'' it \citep{R1}. Agents given opposing personas reproduced 72\% of a human ideological gap in estimates, and 86\% of their reports passed AI methodological review \citep{R3}. A confirmatory prompt flipped Claude Code's verdicts ``from 10\% to 90\% support while leaving its coefficient distribution essentially unchanged'' \citep{R2}, the verdict--estimate gap of Section~\ref{sec:data}. In 640 runs, Claude Opus 4.6 and GPT-5.2-Codex refused direct requests to p-hack, but a prompt reframing specification search as ``uncertainty reporting'' bypassed those guardrails \citep{R5}. A handful of prompt paraphrases can make ``virtually anything'' significant in LLM annotation \citep{R27}, and agents that search against a validation score overfit it \citep{R26}. The integrity tools built for human researchers are therefore more needed with agents, and some transfer. Pre-registering an analysis to run on the next model released would have prevented 17 of 23 and 32 of 44 observed hacks on two tasks \citep{R6}, and a proposed template asks researchers to disclose how many prompt iterations and models they tried \citep{R7}. These are young preprints with different prompts, datasets and models, so the defensible summary is that agents are steerable under some framings, not a rate.

\subsection{The Bridge to Automated Science}

The AI-scientist failures of Section~\ref{sec:science} are the replication crisis, automated. Results chosen after the fact are what one audit calls ``post-hoc selection bias,'' ``akin to training on the test data or p-hacking'' \citep{S10}. Success claimed for a flawed result \citep{S9} is an unchecked verdict, and fabricated results \citep{S15,S16} are what an unverified pipeline lets through. In each case a polished report hides a process no one checked. The remedies are those the human sciences have already developed: a plan and decision rule committed before results exist, the full distribution of defensible analyses, a null check, and an independent re-run by someone other than the analyst.
